\documentclass[12pt,a4paper]{article}

\usepackage[T1]{fontenc}
\usepackage[utf8]{inputenc}
\usepackage[polish]{babel}
\usepackage{lmodern}
\usepackage{geometry}
\usepackage{hyperref}
\usepackage{longtable}
\usepackage{array}
\usepackage{booktabs}
\usepackage{enumitem}
\usepackage{xcolor}
\usepackage{listings}
\usepackage{caption}
\usepackage{graphicx}
\usepackage{float}
\usepackage{tikz}
\usetikzlibrary{shapes.multipart,positioning,arrows.meta}

\geometry{margin=2.5cm}
\setlist[itemize]{leftmargin=1.25cm}
\setlist[enumerate]{leftmargin=1.25cm}
\renewcommand{\arraystretch}{1.25}

\definecolor{codegray}{rgb}{0.95,0.95,0.95}
\definecolor{codeblue}{rgb}{0.0,0.35,0.75}

\lstset{
    backgroundcolor=\color{codegray},
    basicstyle=\ttfamily\small,
    keywordstyle=\color{codeblue}\bfseries,
    frame=single,
    breaklines=true,
    breakatwhitespace=true,
    tabsize=4,
    showstringspaces=false,
    language=PHP
}

\hypersetup{
    colorlinks=true,
    linkcolor=black,
    urlcolor=codeblue,
    pdftitle={Sprawozdanie z pracy projektowej -- Altermatch},
    pdfauthor={Zespół projektowy}
}
\sloppy

% -------------------------------------------------------
% Uzupełnij dane zespołu
% -------------------------------------------------------
\newcommand{\autorA}{Jakub Maliński}
\newcommand{\albumA}{45891}
\newcommand{\autorB}{Michał Kuchar}
\newcommand{\albumB}{44924}
\newcommand{\autorC}{Wojciech Kasprzyk}
\newcommand{\albumC}{44918}
\newcommand{\grupaStudentow}{Collegium Witelona, Legnica \textemdash\ gr.~S2PAM1}
\newcommand{\linkRepo}{https://github.com/Mallinowo/Projekt}
% -------------------------------------------------------

\title{
    \textbf{Sprawozdanie z pracy projektowej}\\[0.5cm]
    \LARGE\textbf{Altermatch}\\[0.3cm]
    \large Aplikacja randkowa dla alternatywnych subkultur
}
\author{
    \begin{tabular}{ll}
        \autorA & (nr albumu: \albumA) \\
        \autorB & (nr albumu: \albumB) \\
        \autorC & (nr albumu: \albumC) \\
    \end{tabular}\\[0.4cm]
    Grupa: \grupaStudentow
}
\date{\today}

\begin{document}

\maketitle
\thispagestyle{empty}

\begin{center}
    \vspace{0.8cm}
    \textbf{Temat projektu:} aplikacja randkowa dla użytkowników związanych z alternatywnymi subkulturami.\\[0.3cm]
    \textbf{Przedmiot:} Programowanie i Projektowanie Systemów Informatycznych 1\\[0.3cm]
    \textbf{Repozytorium:} \url{\linkRepo}
\end{center}

\newpage
\tableofcontents
\newpage

% =======================================================
\section{Wstęp}
% =======================================================

\textbf{Altermatch} to aplikacja randkowa dla użytkowników związanych z alternatywnymi subkulturami: \textit{emo, goth, scene, punk} i \textit{metalhead}. Pomysł wynikał z obserwacji, że mainstreamowe serwisy randkowe projektowane są pod jak najszerszą publiczność -- osoby z silną tożsamością subkulturową często czują się tam nie na miejscu. Zamiast walczyć z algorytmem, który nie wie czym jest \textit{goth}, postanowiliśmy zbudować coś własnego.

Mechanika aplikacji wzoruje się na swipe'ach rozsławionych przez Tindera, ale filtruje kandydatów przez pryzmat subkultury, płci, orientacji, wieku i odległości -- żeby trafiały do nas tylko profile, które faktycznie mogą mieć sens.

Od strony wymagań projektowych Altermatch pokrywa wszystkie 20 zagadnień kwalifikacyjnych: architekturę MVC, ORM, cache, uwierzytelnianie, lokalizację, mailing, asynchroniczne interakcje, konsumpcję i publikację REST API oraz deployment w kontenerach Docker. Stack: \textbf{Laravel 11} (PHP~8.2) na backendzie, \textbf{MySQL~8} i \textbf{Redis} jako warstwa danych i cache, \textbf{Tailwind CSS} i \textbf{Vanilla~JS} na frontendzie, całość opakowana w \textbf{Docker Compose}.

% =======================================================
\section{Opis funkcjonalny systemu}
% =======================================================

\subsection{Charakterystyka ogólna}

Altermatch umożliwia użytkownikom tworzenie profili, przeglądanie kart kandydatów, wyrażanie zainteresowania (like / superlike) lub jego braku (pomiń), a po wzajemnym dopasowaniu -- prowadzenie rozmowy w zintegrowanym czacie. System posiada jedną rolę biznesową: zalogowanego użytkownika z ukończonym onboardingiem.

Główny przepływ biznesowy:
\begin{enumerate}
    \item Użytkownik rejestruje konto z wyborem subkultury, płci i orientacji.
    \item Po rejestracji aplikacja automatycznie geokoduje podane miasto (Open-Meteo Geocoding API) i zapisuje współrzędne potrzebne do filtrowania odległości.
    \item Jednorazowy onboarding zbiera minimum wymaganych danych: zdjęcie profilowe, opis, zainteresowania (min. 5) i opcjonalnie ulubionych artystów.
    \item Moduł \textit{Odkrywaj} wyświetla kandydatów dopasowanych do preferencji użytkownika.
    \item Każda akcja swipe jest wysyłana asynchronicznie do backendu i zapisywana w bazie.
    \item Gdy oboje użytkownicy wykonali akcję like/superlike wobec siebie, tworzony jest match i wysyłane powiadomienia e-mail.
    \item Dopasowane pary mogą rozmawiać w module czatu -- tekstowo, GIF-ami i reakcjami emoji.
\end{enumerate}

\subsection{Rejestracja i logowanie}

Moduł autoryzacji obejmuje:
\begin{itemize}
    \item formularz rejestracji z walidacją: nazwa, e-mail, hasło (min. 8 znaków), wiek (18+), miasto, płeć, orientacja seksualna, subkultura,
    \item automatyczne logowanie po rejestracji i przekierowanie do onboardingu,
    \item formularz logowania z obsługą sesji i opcją \textit{pamiętaj mnie},
    \item wylogowanie: unieważnienie sesji i regeneracja tokenu CSRF,
    \item przełącznik języka interfejsu (PL/EN) -- dostępny bez logowania.
\end{itemize}

Hasła są hashowane przez wbudowany cast modelu Laravel. Formularze są chronione tokenami CSRF generowanymi przez framework.

\subsection{Onboarding użytkownika}

Onboarding jest wymuszony po pierwszej rejestracji -- użytkownik nie ma dostępu do modułu odkrywania dopóki nie uzupełni profilu. Wymagania:
\begin{itemize}
    \item dodanie co najmniej jednego zdjęcia profilowego,
    \item wpisanie opisu profilu (bio),
    \item podanie minimum 5 zainteresowań,
    \item opcjonalne dodanie ulubionych artystów -- ręcznie lub przez Spotify,
    \item potwierdzenie płci i orientacji.
\end{itemize}

Po zakończeniu onboardingu konto otrzymuje flagę \texttt{onboarding\_completed = true}. Middleware w trasach sprawdza tę flagę przy każdym żądaniu do chronionych widoków.

\subsection{Moduł Odkrywaj}

Moduł \texttt{/discover} to centralny ekran aplikacji. Użytkownik widzi stos kart kandydatów. Dostępne akcje:
\begin{itemize}
    \item \textbf{Pomiń (❌)} -- odrzucenie profilu,
    \item \textbf{Like (💚)} -- wyrażenie zainteresowania,
    \item \textbf{Super (⭐)} -- wyróżnione zainteresowanie.
\end{itemize}

Karty obsługuje się przyciskami lub gestem przeciągania (drag). Każda akcja trafia asynchronicznie do \texttt{POST /swipe}. Backend sprawdza, czy druga osoba wcześniej polubiła użytkownika -- jeśli tak, tworzy match.

Kandydaci są filtrowani według preferencji aktywnego użytkownika:
\begin{itemize}
    \item zakres wieku (min–max),
    \item preferowane subkultury (wielokrotny wybór, JSON w kolumnie bazy),
    \item maksymalna odległość w km (Haversine na podstawie \texttt{city\_lat/city\_lng}),
    \item zgodność płci i orientacji (np. orientacja heteroseksualna \textrightarrow\ pokazuj osoby płci przeciwnej),
    \item wykluczenie profili już przesuniętych i już dopasowanych.
\end{itemize}

Lista kandydatów jest cachowana w Redis -- invalidacja następuje po zmianie profilu lub wykonaniu swipe'a.

\subsection{Dopasowania}

Match powstaje po wzajemnym polubieniu. Po jego utworzeniu:
\begin{itemize}
    \item rekord trafia do tabeli \texttt{matches},
    \item interfejs natychmiast pokazuje popup z informacją o nowym dopasowaniu,
    \item obie strony dostają e-mail z powiadomieniem (Laravel Mail),
    \item dopasowane osoby uzyskują dostęp do wspólnego czatu.
\end{itemize}

Polling \texttt{GET /discover/matches/updates} sprawdza co kilka sekund nowe dopasowania i aktualizuje listę w sidebarze.

\subsection{Czat}

Moduł \texttt{/chat} umożliwia komunikację wyłącznie między osobami ze wspólnym matchem. Funkcje:
\begin{itemize}
    \item lista dopasowań w panelu bocznym z podglądem ostatniej wiadomości,
    \item historia wiadomości z paginacją,
    \item wysyłanie wiadomości tekstowych,
    \item wyszukiwanie i wysyłanie GIF-ów (Tenor v2 / GIPHY z automatycznym fallbackiem),
    \item dodawanie reakcji emoji (kliknięcie wiadomości),
    \item oznaczanie wiadomości jako przeczytane przy otwarciu rozmowy,
    \item polling co 3 sekundy -- nowe wiadomości pojawiają się bez przeładowania strony.
\end{itemize}

Backend weryfikuje, że zalogowany użytkownik jest uczestnikiem danego matcha przed zwróceniem jakichkolwiek danych rozmowy.

\subsection{Profil użytkownika}

Moduł profilu pozwala edytować wszystkie dane:
\begin{itemize}
    \item dane podstawowe: nazwa, wiek, miasto, bio,
    \item subkultura, płeć, orientacja,
    \item zdjęcie profilowe i galeria (do 6 zdjęć, upload przez AJAX),
    \item zainteresowania (5--15),
    \item ulubieni artyści (0--10), z opcją synchronizacji ze Spotify,
    \item preferencje odkrywania: zakres wieku, maksymalny dystans, preferowane subkultury.
\end{itemize}

Zmiana danych wpływających na kandydatów (wiek, lokalizacja, subkultura, orientacja) powoduje automatyczne unieważnienie cache discover i publicznego API.

\subsection{Integracja ze Spotify}

Spotify jest opcjonalną integracją -- aplikacja działa w pełni bez niej. Obsługiwany przepływ:
\begin{enumerate}
    \item Użytkownik klika \textit{Połącz z Spotify} w profilu.
    \item Backend przekierowuje na OAuth2 Spotify z scope \texttt{user-top-read}.
    \item Po autoryzacji Spotify wraca na \texttt{/spotify/callback} z kodem.
    \item Backend wymienia kod na token, zapisuje go i oferuje listę ulubionych artystów.
    \item Użytkownik może jednym kliknięciem dodać artystów do profilu.
\end{enumerate}

\subsection{Lokalizacja językowa}

Aplikacja obsługuje dwa języki: polski (domyślny) i angielski. Tłumaczenia znajdują się w \texttt{resources/lang/pl} i \texttt{resources/lang/en}. Wybór języka zapisywany jest w sesji oraz w kolumnie \texttt{locale} użytkownika w bazie danych. Przełącznik działa również dla niezalogowanych gości.

\subsection{Publiczne REST API}

Projekt publikuje API w prefiksie \texttt{/api/v1}:
\begin{itemize}
    \item \texttt{GET /api/v1/profiles} -- paginowana lista profili (cacheowana),
    \item \texttt{GET /api/v1/profiles/\{id\}} -- szczegóły profilu po ID,
    \item \texttt{GET /api/v1/subcultures} -- słownik subkultur z ikonami.
\end{itemize}

Endpointy są publiczne (nie wymagają tokenu) i odpowiadają JSON-em.

% =======================================================
\section{Opis technologiczny}
% =======================================================

\subsection{Architektura aplikacji}

System jest monolityczną aplikacją MVC opartą na konwencjach Laravel. Warstwy:

\begin{itemize}
    \item \textbf{Model} -- klasy Eloquent: \texttt{User}, \texttt{Photo}, \texttt{Interest}, \texttt{Artist}, \texttt{Swipe}, \texttt{UserMatch}, \texttt{Message}, \texttt{MessageReaction},
    \item \textbf{View} -- szablony Blade renderowane po stronie serwera,
    \item \textbf{Controller} -- klasy obsługujące żądania HTTP; logika biznesowa jest częściowo wydzielona do klas \textbf{Service},
    \item \textbf{Service} -- \texttt{SpotifyService}, \texttt{GifService}, \texttt{CityLocationService} -- integracje z zewnętrznymi API.
\end{itemize}

Uproszczony diagram przepływu żądania:

\begin{lstlisting}[language={}]
Browser -> Nginx -> PHP-FPM/Laravel
  |
  +-- Middleware: auth, onboarding_completed, locale
  +-- Router: web.php / api.php
  +-- Controller -> Model (Eloquent) -> MySQL
  |              -> Cache (Redis)
  |              -> Service (Spotify / GIF / Geocoding)
  +-- View: Blade template -> HTML -> Browser
\end{lstlisting}

\subsection{Backend}

Backend działa na \textbf{PHP 8.2} z \textbf{Laravel 11}. Najważniejsze kontrolery:

\begin{itemize}
    \item \texttt{AuthController} -- rejestracja, logowanie, wylogowanie, geokodowanie miasta, zmiana locale,
    \item \texttt{OnboardingController} -- walidacja i zapis danych onboardingu, upload zdjęcia,
    \item \texttt{DiscoverController} -- zapytanie do bazy z wielowarstwowym filtrowaniem, obsługa cache, lista nowych dopasowań (polling),
    \item \texttt{SwipeController} -- atomowy zapis swipe'a, detekcja wzajemnego lajku, tworzenie matcha, wysyłka maili,
    \item \texttt{ChatController} -- wczytanie wiadomości, wysyłka, wyszukiwanie GIF-ów, zapis reakcji, oznaczanie przeczytanych,
    \item \texttt{ProfileController} -- edycja danych profilu, upload do 6 zdjęć (Intervention Image), unieważnienie cache,
    \item \texttt{SpotifyController} -- OAuth2 flow, search artystów, pobieranie top artists, odłączenie konta,
    \item \texttt{ApiController} -- publiczne REST API z paginacją i cache.
\end{itemize}

Walidacja danych wejściowych opiera się na mechanizmie Laravel Request Validation. Każdy kontroler zwraca albo widok Blade, albo odpowiedź JSON.

\subsection{Baza danych}

Warstwa danych korzysta z \textbf{MySQL 8}. Struktura jest zarządzana przez migracje Laravel. Główne tabele:

\begin{longtable}{|p{3.8cm}|p{9.4cm}|}
\hline
\textbf{Tabela} & \textbf{Zawartość} \\
\hline
\texttt{users} & Dane konta, profilu, geolokalizacja, preferencje odkrywania, locale, flaga onboardingu \\
\hline
\texttt{photos} & Zdjęcia użytkownika z kolejnością wyświetlania \\
\hline
\texttt{interests} & Zainteresowania użytkownika (relacja 1:N do \texttt{users}) \\
\hline
\texttt{artists} & Ulubieni artyści, opcjonalnie powiązani ze Spotify \\
\hline
\texttt{swipes} & Historia decyzji: kto, kogo, jaką akcją (like/dislike/superlike) \\
\hline
\texttt{matches} & Pary wzajemnie dopasowanych użytkowników \\
\hline
\texttt{messages} & Wiadomości czatu: typ (text/gif), treść, gif\_url, read\_at \\
\hline
\texttt{message\_reactions} & Reakcje emoji na wiadomości \\
\hline
\end{longtable}

Relacje Eloquent:
\begin{itemize}
    \item \texttt{User} \textrightarrow\ \texttt{hasMany} Photo, Interest, Artist, Swipe,
    \item \texttt{User} \textrightarrow\ \texttt{belongsToMany} User (przez \texttt{matches}),
    \item \texttt{UserMatch} \textrightarrow\ \texttt{hasMany} Message,
    \item \texttt{Message} \textrightarrow\ \texttt{hasMany} MessageReaction.
\end{itemize}

\begin{figure}[H]
\centering
\resizebox{\textwidth}{!}{%
\begin{tikzpicture}[
    node distance=1.8cm and 3.5cm,
    entity/.style={
        rectangle split, rectangle split parts=2,
        draw=black, thick,
        rectangle split part fill={blue!18, blue!4},
        font=\scriptsize\ttfamily,
        align=left, inner sep=3pt
    },
    arr/.style={-{Stealth[length=5pt,width=4pt]}, thick},
    darr/.style={-{Stealth[length=5pt,width=4pt]}, thick, dashed, draw=gray},
    lbl/.style={font=\tiny\sffamily, fill=white, inner sep=1pt, midway}
]

\node[entity] (users) {%
    \textbf{users}%
    \nodepart{two}%
    \underline{id}\ name\ email\ password\\%
    age\ gender\ orientation\ subculture\\%
    city\ city\_lat\ city\_lng\\%
    bio\ avatar\ locale\\%
    onboarding\_completed\\%
    pref\_age\_min\ pref\_age\_max\\%
    pref\_max\_distance\_km\\%
    pref\_subcultures (JSON)%
};

\node[entity, above left=1.5cm and 5.0cm of users] (photos) {%
    \textbf{photos}%
    \nodepart{two}%
    \underline{id}\ \textit{user\_id}\ path\ order%
};

\node[entity, above=3.2cm of users] (interests) {%
    \textbf{interests}%
    \nodepart{two}%
    \underline{id}\ \textit{user\_id}\ name%
};

\node[entity, above right=1.5cm and 5.0cm of users] (artists) {%
    \textbf{artists}%
    \nodepart{two}%
    \underline{id}\ \textit{user\_id}\ name%
};

\node[entity, below left=2.5cm and 3.2cm of users] (swipes) {%
    \textbf{swipes}%
    \nodepart{two}%
    \underline{id}\\%
    \textit{swiper\_id}\ \textit{swiped\_id}\\%
    direction%
};

\node[entity, right=5.5cm of users] (matches) {%
    \textbf{matches}%
    \nodepart{two}%
    \underline{id}\\%
    \textit{user1\_id}\ \textit{user2\_id}%
};

\node[entity, below=2.5cm of matches] (messages) {%
    \textbf{messages}%
    \nodepart{two}%
    \underline{id}\ \textit{match\_id}\ \textit{sender\_id}\\%
    type\ body\ gif\_url\ read\_at%
};

\node[entity, below=2.0cm of messages] (reactions) {%
    \textbf{message\_reactions}%
    \nodepart{two}%
    \underline{id}\ \textit{message\_id}\ \textit{user\_id}\\%
    emoji%
};

\draw[arr] (photos)    -- node[lbl,above,sloped]{N:1} (users);
\draw[arr] (interests) -- node[lbl,left]{N:1}         (users);
\draw[arr] (artists)   -- node[lbl,above,sloped]{N:1} (users);
\draw[arr] (swipes)    -- node[lbl,above,sloped]{N:1} (users);
\draw[arr] (matches)   -- node[lbl,above]{N:M}        (users);
\draw[arr] (messages)  -- node[lbl,right]{N:1}        (matches);
\draw[arr] (reactions) -- node[lbl,right]{N:1}        (messages);
\draw[darr] (messages.west) to[out=180,in=270]
    node[lbl,below,near end]{sender\_id} (users.south);
\draw[darr] (reactions.west) to[out=180,in=250,looseness=0.8]
    node[lbl,below,near start]{user\_id} (users.south west);

\end{tikzpicture}%
}
\caption{Diagram ERD -- schemat bazy danych aplikacji Altermatch. \underline{Podkreślenie} -- klucz główny (PK), \textit{kursywa} -- klucz obcy (FK). Linie przerywane oznaczają dodatkowe relacje FK (\texttt{sender\_id}, \texttt{user\_id}).}
\label{fig:erd}
\end{figure}

\subsection{Cache i sesje}

Projekt używa \textbf{Redis} (przez bibliotekę Predis) jako backendu zarówno dla cache aplikacji, jak i dla sesji. Strategie cache:

\begin{itemize}
    \item \textbf{Lista kandydatów discover} -- klucz zależny od ID użytkownika i wersji cache; wersja jest inkrementowana po każdym swipe'ie lub zmianie profilu,
    \item \textbf{Endpointy API} -- wyniki \texttt{/api/v1/profiles} i \texttt{/api/v1/subcultures} cachowane przez 10 minut,
    \item \textbf{Token Spotify client credentials} -- cache na czas ważności tokenu,
    \item \textbf{Wyniki wyszukiwania GIF-ów} -- cache na 5 minut (zmniejsza liczbę zapytań do Tenor/GIPHY),
    \item \textbf{Wyniki geokodowania} -- cache bezterminowy (miasto do współrzędnych zmienia się rzadko).
\end{itemize}

\subsection{Frontend}

Warstwa widoku:
\begin{itemize}
    \item \textbf{Blade} -- szablony HTML renderowane po stronie serwera; layouty \texttt{app} (zalogowani) i \texttt{auth} (gość),
    \item \textbf{Tailwind CSS} -- framework utility-first; responsywność przez klasy \texttt{sm:}, \texttt{md:}, \texttt{lg:},
    \item \textbf{Vite} -- bundler zasobów frontendowych (CSS, JS), hot reload w trybie dev,
    \item \textbf{Vanilla JavaScript} -- bez zewnętrznych frameworków JS.
\end{itemize}

Kluczowe skrypty JavaScript:
\begin{itemize}
    \item \textbf{public/js/discover.js} -- animacje kart (drag + CSS transform), obsługa przycisków swipe, popup nowego matcha, polling listy dopasowań co 5 sekund,
    \item \textbf{public/js/chat.js} -- otwieranie rozmów przez Fetch API, wysyłanie wiadomości, wyszukiwanie GIF-ów, zapis reakcji emoji, polling co 3 sekundy.
\end{itemize}

\subsection{Integracje zewnętrzne}

\begin{longtable}{|p{3.5cm}|p{9.7cm}|}
\hline
\textbf{Serwis} & \textbf{Sposób użycia} \\
\hline
\textbf{Spotify API} & OAuth2 (authorization code flow), endpoint \texttt{/v1/search}, \texttt{/v1/me/top/artists}. Token client credentials cachowany w Redis. \\
\hline
\textbf{Tenor v2 API} & Wyszukiwanie GIF-ów w czacie; fallback na Tenor v1 i GIPHY. Wyniki cachowane. \\
\hline
\textbf{GIPHY API} & Fallback jeśli Tenor jest niedostępny lub bez klucza. \\
\hline
\textbf{Open-Meteo Geocoding} & Zamiana nazwy miasta na \texttt{latitude/longitude} przy rejestracji. Wyniki cachowane. \\
\hline
\end{longtable}

Wszystkie integracje są opakowane blokami \texttt{try/catch}. Błędy są logowane przez Laravel Log -- aplikacja nie przerywa działania gdy zewnętrzne API jest niedostępne.

\subsection{Mailing}

Laravel Mail z driverem SMTP. W środowisku deweloperskim wiadomości przechwytuje \textbf{Mailpit} (port 8025). Wysyłane wiadomości:
\begin{itemize}
    \item \textbf{Welcome} -- po zakończeniu rejestracji,
    \item \textbf{Match} -- gdy dwoje użytkowników wzajemnie się polubilo; wysyłany do obojga.
\end{itemize}

Widoki e-maili są szablonami Blade w katalogu \texttt{resources/views/emails}.

\subsection{Logowanie zdarzeń}

Laravel Log (Monolog) zapisuje zdarzenia do pliku \texttt{storage/logs/laravel.log}. Logowane są m.in.:
\begin{itemize}
    \item rejestracja nowego użytkownika (ID, e-mail),
    \item błędy geokodowania miasta,
    \item tworzenie matcha,
    \item błędy wysyłki poczty,
    \item błędy integracji Spotify i GIF API.
\end{itemize}

\subsection{Konteneryzacja i wdrożenie}

Środowisko opisano w \texttt{docker-compose.yml}. Usługi:

\begin{longtable}{|p{2.5cm}|p{4.0cm}|p{6.7cm}|}
\hline
\textbf{Usługa} & \textbf{Technologia} & \textbf{Rola} \\
\hline
\texttt{app} & PHP 8.2-FPM, Composer, Node.js & Aplikacja Laravel, kompilacja assetów \\
\hline
\texttt{nginx} & Nginx & Serwer HTTP, routing, przyjazne URL \\
\hline
\texttt{mysql} & MySQL 8 & Relacyjna baza danych \\
\hline
\texttt{redis} & Redis 7 & Cache, sesje \\
\hline
\texttt{mailpit} & Mailpit & Lokalne przechwytywanie e-maili (dev) \\
\hline
\end{longtable}

Dockerfile instaluje zależności PHP (Composer) i Node.js (npm), buduje assety frontendowe przez Vite, ustawia uprawnienia katalogów \texttt{storage} i \texttt{bootstrap/cache}.

% =======================================================
\section{Wyszczególnione wdrożone zagadnienia kwalifikacyjne}
% =======================================================

\begin{longtable}{|>{\centering\arraybackslash}p{0.7cm}|p{3.4cm}|p{8.0cm}|p{1.6cm}|}
\hline
\textbf{\#} & \textbf{Zagadnienie} & \textbf{Realizacja w projekcie} & \textbf{Status} \\
\hline
1 & Framework MVC & Laravel 11 -- kontrolery, modele Eloquent i widoki Blade. Podział na warstwy Model, View, Controller i Service jest konsekwentnie zachowany. & Wdrożone \\
\hline
2 & Framework CSS & Tailwind CSS z konfiguracją Vite i PostCSS. Responsywne klasy \texttt{sm:}, \texttt{md:}, \texttt{lg:} w widokach. & Wdrożone \\
\hline
3 & Baza danych & MySQL 8. Struktura zarządzana przez migracje. Seedery tworzą konta demo i dane testowe. & Wdrożone \\
\hline
4 & Cache & Redis przez Laravel Cache i bibliotekę Predis. Cachowane są m.in.: lista discover, endpointy API, tokeny Spotify, wyniki GIF, geokodowanie. & Wdrożone \\
\hline
5 & Dependency manager & Composer dla PHP (Laravel, Sanctum, Predis, Intervention Image). npm + Vite dla zasobów frontendowych (Tailwind, Alpine). & Wdrożone \\
\hline
6 & HTML & Semantyczne szablony Blade: layouty, widoki auth, onboarding, discover, chat, profile, e-maile. & Wdrożone \\
\hline
7 & CSS & Tailwind CSS oraz własne style komponentów, animacji kart i modalu matcha. & Wdrożone \\
\hline
8 & JavaScript & Vanilla JS: animacje drag-to-swipe, Fetch API dla akcji swipe/chat/upload, polling, dynamiczny UI, GIF-y, reakcje emoji. & Wdrożone \\
\hline
9 & Routing & Laravel Router -- grupy middleware, named routes, zasoby RESTful. Nginx obsługuje przyjazne URL (pretty URLs) i przekazuje żądania do PHP-FPM. & Wdrożone \\
\hline
10 & ORM & Eloquent ORM -- modele, relacje (\texttt{hasMany}, \texttt{belongsTo}, \texttt{belongsToMany}), eager loading, scopy zapytań. & Wdrożone \\
\hline
11 & Uwierzytelnianie & Rejestracja i logowanie przez Laravel Auth. Sesje chronione przez middleware \texttt{auth}. Sanctum przygotowany dla API tokenowego. CSRF na wszystkich formularzach. & Wdrożone \\
\hline
12 & Lokalizacja & Laravel i18n -- pliki tłumaczeń PL/EN w \texttt{resources/lang}. Middleware locale ustawia język na podstawie sesji lub kolumny \texttt{locale} użytkownika. & Wdrożone \\
\hline
13 & Mailing & Laravel Mail -- Mailable \textit{Welcome} i \textit{Match}. Widoki e-maili w Blade. W środowisku deweloperskim Mailpit przechwytuje wiadomości pod \texttt{localhost:8025}. & Wdrożone \\
\hline
14 & Formularze & Formularze Blade z tokenem CSRF, walidacja serwerowa przez Laravel Request Validation, zwrotne komunikaty błędów przy polach. & Wdrożone \\
\hline
15 & Asynchroniczność & Fetch API -- akcje swipe, wysyłanie i odbieranie wiadomości, upload zdjęć, wyszukiwanie GIF-ów. Polling discover i czatu -- aktualizacja UI bez przeładowania strony. & Wdrożone \\
\hline
16 & Konsumpcja API & Spotify API (OAuth2, artyści), Tenor v2/GIPHY (GIF-y), Open-Meteo Geocoding (geolokalizacja miast). Klient HTTP: Laravel Http/Guzzle. & Wdrożone \\
\hline
17 & Publikacja API & REST API \texttt{/api/v1/profiles} (lista + paginacja), \texttt{/api/v1/profiles/\{id\}}, \texttt{/api/v1/subcultures}. Odpowiedzi JSON, cachowane w Redis. & Wdrożone \\
\hline
18 & RWD & Responsywne klasy Tailwind. Interfejs dostosowany do różnych szerokości: mobile (pojedyncza kolumna), tablet i desktop (sidebar + główna treść). & Wdrożone \\
\hline
19 & Logger & Laravel Log (Monolog) -- rejestracja zdarzeń aplikacji: rejestracja użytkownika, tworzenie matcha, błędy API, błędy maili, zdarzenia Spotify. & Wdrożone \\
\hline
20 & Deployment & Dockerfile (PHP-FPM, Composer, Node.js, Vite). Docker Compose z usługami: \texttt{app}, \texttt{nginx}, \texttt{mysql}, \texttt{redis}, \texttt{mailpit}. & Wdrożone \\
\hline
\caption{Zestawienie 20 zagadnień kwalifikacyjnych}
\end{longtable}

Wszystkie 20 zagadnień kwalifikacyjnych zostało zrealizowanych, co kwalifikuje projekt do oceny \textbf{5.0} według skali opisanej w wymaganiach.

% =======================================================
\section{Instrukcja lokalnego i zdalnego uruchomienia}
% =======================================================

\subsection{Wymagania lokalne}

\begin{itemize}
    \item Docker (wersja 24+),
    \item Docker Compose (wbudowany w Docker Desktop lub jako plugin),
    \item wolne porty: \texttt{8080} (aplikacja), \texttt{8025} (Mailpit), \texttt{3307} (MySQL),
    \item połączenie z internetem przy pierwszym budowaniu obrazu.
\end{itemize}

\subsection{Uruchomienie lokalne}

\begin{enumerate}

    \item Sklonować repozytorium:
\begin{lstlisting}[language=bash]
git clone <adres_repozytorium>
cd altermatch
\end{lstlisting}

    \item Przygotować plik środowiskowy:
\begin{lstlisting}[language=bash]
cp .env.example .env
\end{lstlisting}

    \item Zbudować i uruchomić kontenery:
\begin{lstlisting}[language=bash]
docker compose up -d --build
\end{lstlisting}

    \item Wygenerować klucz aplikacji:
\begin{lstlisting}[language=bash]
docker compose exec app php artisan key:generate
\end{lstlisting}

    \item Uruchomić migracje i seedery (tworzy konta demo):
\begin{lstlisting}[language=bash]
docker compose exec app php artisan migrate --seed
\end{lstlisting}

    \item Utworzyć symlink do plików publicznych w \texttt{storage}:
\begin{lstlisting}[language=bash]
docker compose exec app php artisan storage:link
\end{lstlisting}

\end{enumerate}

\subsection{Adresy lokalnych usług}

\begin{itemize}
    \item Aplikacja: \url{http://localhost:8080}
    \item Mailpit (podgląd e-maili): \url{http://localhost:8025}
    \item MySQL: \texttt{localhost:3307} (dane w \texttt{.env})
\end{itemize}

\subsection{Konta demonstracyjne}

Hasło dla wszystkich kont demo: \texttt{demo123}.

\begin{center}
\begin{tabular}{lll}
\toprule
\textbf{E-mail} & \textbf{Imię} & \textbf{Subkultura} \\
\midrule
\texttt{moon@demo.pl} & Klaudia & emo \\
\texttt{void@demo.pl} & Krystian & goth \\
\texttt{goth@demo.pl} & Julka & scene \\
\texttt{rain@demo.pl} & Oliwier & punk \\
\texttt{crypt@demo.pl} & Luciusz & metalhead \\
\bottomrule
\end{tabular}
\end{center}

\subsection{Kluczowe zmienne środowiskowe}

\begin{longtable}{|p{5.5cm}|p{7.7cm}|}
\hline
\textbf{Zmienna} & \textbf{Opis} \\
\hline
\texttt{APP\_URL} & Adres aplikacji; \texttt{http://localhost:8080} lokalnie \\
\hline
\texttt{DB\_HOST}, \texttt{DB\_DATABASE} & Adres i nazwa bazy MySQL \\
\hline
\texttt{CACHE\_STORE=redis} & Backend cache aplikacji \\
\hline
\texttt{SESSION\_DRIVER=redis} & Backend sesji \\
\hline
\texttt{MAIL\_HOST=mailpit} & Serwer pocztowy (dev) \\
\hline
\texttt{TENOR\_API\_KEY} & Klucz API Tenor (opcjonalny) \\
\hline
\texttt{GIPHY\_API\_KEY} & Klucz API GIPHY (opcjonalny) \\
\hline
\texttt{SPOTIFY\_CLIENT\_ID} & Spotify OAuth2 Client ID (opcjonalny) \\
\hline
\texttt{SPOTIFY\_CLIENT\_SECRET} & Spotify OAuth2 Secret (opcjonalny) \\
\hline
\texttt{SPOTIFY\_REDIRECT\_URI} & Spotify callback URI (opcjonalny) \\
\hline
\end{longtable}

Klucze API Spotify, Tenor i GIPHY są opcjonalne -- aplikacja działa bez nich; wyłączone są tylko odpowiednie funkcje.

\subsection{Uruchomienie zdalne (VPS)}

Wdrożenie produkcyjne na serwerze Linux z Dockerem:
\begin{enumerate}
    \item Zainstalować Docker i Docker Compose.
    \item Skonfigurować DNS domeny na IP serwera.
    \item Sklonować repozytorium na serwer.
    \item Przygotować plik \texttt{.env} z ustawieniami produkcyjnymi:
    \begin{itemize}
        \item \texttt{APP\_ENV=production}, \texttt{APP\_DEBUG=false},
        \item wygenerowany \texttt{APP\_KEY},
        \item silne hasło bazy MySQL,
        \item prawidłowy \texttt{APP\_URL} z domeną.
    \end{itemize}
    \item Uruchomić kontenery: \texttt{docker compose up -d --build}
    \item Wykonać migracje: \texttt{docker compose exec app php artisan migrate --force}
    \item Utworzyć symlink: \texttt{docker compose exec app php artisan storage:link}
    \item Skonfigurować reverse proxy (Nginx/Caddy) i certyfikat HTTPS (Let's Encrypt).
    \item Przygotować backup bazy MySQL i katalogu \texttt{storage/app/public}.
\end{enumerate}

\subsection{Generowanie PDF ze sprawozdania}

\begin{lstlisting}[language=bash]
pdflatex sprawozdanie.tex
pdflatex sprawozdanie.tex   # drugi przebieg: spis tresci i odnosniki
\end{lstlisting}

Wymagane: TeX Live lub MiKTeX.

% =======================================================
\section{Podział pracy w zespole}
% =======================================================

% Uzupełnij rzeczywisty podział pracy
\begin{longtable}{|p{4.5cm}|p{8.7cm}|}
\hline
\textbf{Członek zespołu} & \textbf{Zakres odpowiedzialności} \\
\hline
\autorA & Backend: modele Eloquent i migracje, \texttt{DiscoverController} (filtrowanie kandydatów, algorytm Haversine, cache Redis), \texttt{SwipeController} (logika matchy), publiczne REST API, konfiguracja Redis \\
\hline
\autorB & Frontend: szablony Blade, Tailwind CSS, responsywność RWD, JavaScript (\texttt{discover.js} -- animacje drag-to-swipe, popup matcha; \texttt{chat.js} -- polling, GIF-y, reakcje emoji), widoki onboardingu i profilu \\
\hline
\autorC & Integracje i infrastruktura: Spotify OAuth2, Tenor\,v2/GIPHY z fallbackiem, geokodowanie Open-Meteo, Docker Compose, konfiguracja Nginx, mailing (Laravel Mail + Mailpit), \texttt{AuthController}, \texttt{ProfileController} \\
\hline
\end{longtable}

% =======================================================
\section{Wnioski projektowe}
% =======================================================

\subsection{Podsumowanie realizacji}

Projekt Altermatch zrealizowano jako kompletną aplikację webową pokrywającą wszystkie 20 zagadnień kwalifikacyjnych. System przeszedł przez pełen cykl wytwórczy: projektowanie bazy danych i relacji, implementację backendu Laravel, budowę interfejsu Blade/Tailwind, dodanie dynamicznych interakcji JavaScript, integrację z zewnętrznymi API oraz konteneryzację przez Docker Compose.

Architektura monolityczna MVC okazała się właściwym wyborem dla tego zakresu funkcjonalności -- pozwoliła utrzymać spójność kodu przy jednoczesnym zachowaniu podziału odpowiedzialności (kontrolery, modele, serwisy, widoki). Laravel dostarczył gotowe mechanizmy dla najbardziej złożonych wymagań: uwierzytelniania, walidacji, ORM, cache, maili i routingu.

\subsection{Napotkane trudności}

\begin{itemize}
    \item \textbf{Spójność stanu discover/czat}: Utrzymanie synchronizacji między stanem frontend (JS) a backendem bez WebSocketów wymagało przemyślanej strategii pollingu i inwalidacji cache. Rozwiązano to przez wersjonowanie kluczy cache i polling z krótkim interwałem.

    \item \textbf{Filtrowanie z odległością}: Obliczanie odległości Haversine bezpośrednio w zapytaniu SQL na dużej tabeli może być kosztowne. Zoptymalizowano to przez caching wyników i filtr wstępny po kolumnach lat/lng przed obliczeniem dokładnej odległości.

    \item \textbf{Inwalidacja cache}: Cache discover musi być unieważniany po swipe'ie, zmianie profilu, ale nie po każdym żądaniu. Zastosowano wersjonowanie klucza cache per użytkownik.

    \item \textbf{Bezpieczeństwo czatu}: Dostęp do wiadomości musi być ograniczony do uczestników danego matcha. Każde żądanie do \texttt{ChatController} weryfikuje, że zalogowany użytkownik jest jedną ze stron dopasowania.

    \item \textbf{Fallback integracji zewnętrznych}: Aplikacja musi działać nawet gdy Tenor, GIPHY lub Spotify są niedostępne lub nie skonfigurowane. Zaimplementowano wielowarstwowy fallback (Tenor v2 → Tenor v1 → GIPHY → brak GIF-ów) i opakowanie wyjątków w bloki try/catch.
\end{itemize}

\subsection{Najważniejsze rezultaty}

\begin{itemize}
    \item Wdrożono pełen przepływ użytkownika: rejestracja → onboarding → discover → match → czat.
    \item Zaimplementowano inteligentne filtrowanie kandydatów uwzględniające wiek, odległość, subkulturę i zgodność płci/orientacji.
    \item Zintegrowano trzy zewnętrzne API z fallbackiem i cachowaniem.
    \item Opublikowano publiczne REST API z paginacją i cachowaniem.
    \item Skonteneryzowano całe środowisko (app, baza, cache, poczta) w Docker Compose.
    \item Wszystkie 20 wymaganych zagadnień kwalifikacyjnych zostało wdrożonych.
\end{itemize}

\subsection{Możliwe kierunki dalszego rozwoju}

\begin{itemize}
    \item Zastąpienie pollingu WebSocketami (np. Laravel Reverb / Soketi) dla komunikacji w czasie rzeczywistym.
    \item Dodanie testów automatycznych: jednostkowych, integracyjnych i end-to-end (PHPUnit, Pest, Laravel Dusk).
    \item Wdrożenie CI/CD (GitHub Actions) -- automatyczne testy i deploy po każdym push.
    \item Panel administracyjny z narzędziami moderacji treści.
    \item Rozbudowa REST API o autoryzację przez tokeny API (Sanctum) i endpointy zapisu.
    \item Wdrożenie kolejek (Laravel Queue + Redis) dla wysyłki maili w tle.
    \item Monitoring i alerty (Sentry, Prometheus, Grafana) dla środowiska produkcyjnego.
\end{itemize}

\subsection{Wniosek końcowy}

Altermatch był dla nas dobrym projektem -- wystarczająco złożonym, żeby każde zagadnienie kwalifikacyjne miało sens w kontekście realnej funkcjonalności, a nie \textit{feature for the sake of feature}. Praca nad spójnością stanu między frontendem a backendem bez WebSocketów, debugowanie inwalidacji cache czy przepływ OAuth2 Spotify nauczyły nas więcej niż niejedne ćwiczenia laboratoryjne.

Projekt spełnia wszystkie 20 wymagań kwalifikacyjnych, jest kompletny i działa w środowisku Docker. Zastosowany stack (Laravel~11, MySQL, Redis, Tailwind CSS, Docker) jest reprezentatywny dla współczesnych aplikacji webowych i stanowi solidną bazę do dalszego rozwoju.

\end{document}
